\chapter{Допуск общего configuration space конечных endpoint-геометрий}
\label{ch:v9-endpoint-finite-geometry-configuration-space}

\section{Три фазы как ортогональные компоненты}

Предыдущий no-go требует расширить не одну fibre, а пространство допустимых
конечных геометрий. Зафиксированное ранее endpoint-меню имеет вложения
\begin{equation}
 H_{21}\subset H_{23}=H_{21}\oplus H_{\rm n}^{2}
 \subset H_{24}=H_{23}\oplus H_{\rm t}^{1},
 \qquad (\dim H_{21},\dim H_{23},\dim H_{24})=(21,23,24).
 \label{eq:v9-config-space-phase-chain}
\end{equation}
Рассмотрим конечное base-space и его коммутативную phase algebra
\begin{equation}
 X_F=\{21,23,24\},
 \qquad \mathcal Z_F=C(X_F)\simeq\mathbb C^3.
 \label{eq:v9-config-space-phase-algebra}
\end{equation}
Минимальное faithful orthogonal представление всех трёх fibres есть
\begin{equation}
 H_\Sigma=H_{21}\oplus H_{23}\oplus H_{24},
 \qquad \dim_{\mathbb C}H_\Sigma=21+23+24=68.
 \label{eq:v9-config-space-total-carrier}
\end{equation}
Центральные phase-projectors удовлетворяют
\begin{equation}
 Q_iQ_j=\delta_{ij}Q_i,
 \qquad Q_{21}+Q_{23}+Q_{24}=I_{68},
 \qquad (\rank Q_{21},\rank Q_{23},\rank Q_{24})=(21,23,24).
 \label{eq:v9-config-space-central-projectors}
\end{equation}
Таким образом, включение $H_{24}$ означает наличие кандидата, но не выбор
физической фазы.

\section{Канонические рёбра и общий Dirac-блок}

Пусть $V:H_{21}\to H_{23}$ и $W:H_{23}\to H_{24}$ --- канонические
изометрии вложений. Их defect-projectors имеют точные ранги
\begin{equation}
 V^*V=I_{21},\qquad VV^*=I_{23}-P_{\rm n},
 \qquad \rank P_{\rm n}=2,
 \label{eq:v9-config-space-neutral-isometry}
\end{equation}
\begin{equation}
 W^*W=I_{23},\qquad WW^*=I_{24}-P_{\rm t},
 \qquad \rank P_{\rm t}=1.
 \label{eq:v9-config-space-triplet-isometry}
\end{equation}
На $H_\Sigma$ определён ближайший соседний self-adjoint operator
\begin{equation}
 D_\Sigma=
 \begin{pmatrix}
 0&t_{\rm n}V^*&0\\
 t_{\rm n}V&0&t_{\rm t}W^*\\
 0&t_{\rm t}W&0
 \end{pmatrix},
 \qquad t_{\rm n},t_{\rm t}\in\mathbb R\setminus\{0\}.
 \label{eq:v9-config-space-dirac-block}
\end{equation}
Для unit weights точный аудит даёт
\begin{equation}
 D_\Sigma^*=D_\Sigma,
 \qquad \rank D_\Sigma=46,
 \qquad \dim\ker D_\Sigma=22.
 \label{eq:v9-config-space-dirac-rank}
\end{equation}
Это конечный quiver-type Dirac на уже заданном representation graph
\cite{v9-config-space-krajewski,v9-config-space-quiver-action}.

\section{Связность без выбора финальной вершины}

Underlying path graph имеет Laplacian
\begin{equation}
 L_F=\begin{pmatrix}1&-1&0\\-1&2&-1\\0&-1&1\end{pmatrix}.
 \label{eq:v9-config-space-laplacian}
\end{equation}
Его spectrum и kernel равны
\begin{equation}
 \Spec L_F=\{0,1,3\},
 \qquad \rank L_F=2,
 \qquad \ker L_F=\mathbb C(1,1,1)^T.
 \label{eq:v9-config-space-laplacian-spectrum}
\end{equation}
Поэтому чистая graph-кинетика
\begin{equation}
 S_{\rm graph}(z)=z^*L_Fz\geq0,
 \qquad \operatorname*{argmin}S_{\rm graph}=\mathbb C(1,1,1)^T
 \label{eq:v9-config-space-graph-minimum}
\end{equation}
выбирает constant mode, а не одну из фаз $21,23,24$. Никакой target-loaded
вершинный score в configuration space не вставлен.

\section{Reachability obstruction канонических вложений}

Три копии старого модуля образуют подпространство
\begin{equation}
 H_{\rm old}^{\Sigma}
 =H_{21}\oplus V H_{21}\oplus WV H_{21},
 \qquad \dim H_{\rm old}^{\Sigma}=63.
 \label{eq:v9-config-space-old-subbundle}
\end{equation}
Оно редуцирует канонический Dirac-блок:
\begin{equation}
 [D_\Sigma,P_{\rm old}^{\Sigma}]=0.
 \label{eq:v9-config-space-old-invariance}
\end{equation}
Ортогональный остаток содержит две neutral-линии в fibre $23$, их две
копии в fibre $24$ и новую charged triplet-линию:
\begin{equation}
 \dim(H_{\rm old}^{\Sigma})^\perp=68-63=5=2+2+1.
 \label{eq:v9-config-space-unreachable-defect}
\end{equation}
Следовательно, для любого polynomial evolution
\begin{equation}
 P_{\rm new}^{\Sigma}p(D_\Sigma)Q_{21}=0
 \quad\text{для всех }p,
 \qquad N_{\rm created\ endpoint\ lines}=0/3.
 \label{eq:v9-config-space-reachability-no-go}
\end{equation}
Прямая сумма и connected phase graph связывают названия фаз, но
канонические inclusions переносят только уже существующие состояния.

Есть и family-дефект фиксированного triplet-вложения. На charged
$\mathbb C^3$ projector старых двух координат не инвариантен:
\begin{equation}
 P_{\rm old,t}=\operatorname{diag}(0,1,1),
 \qquad \rank[J_{12},P_{\rm old,t}]=2.
 \label{eq:v9-config-space-family-edge-defect}
\end{equation}
Поэтому complete transition-edge должен сам нести family-vector type, а не
быть фиксированной scalar-изометрией.

\section{Итог допуска}

Общее конечное configuration space построено без выбора target vertex, но
его канонический Dirac не является оператором рождения:
\begin{equation}
 \begin{aligned}
 N_{\rm configuration\ space\ architecture}&=9/9,\\
 N_{\rm endpoint\ creation\ reachability}&=0/3,\\
 N_{\rm physical\ phase\ selector}&=0/1,\\
 N_{\rm physical\ four\ slot\ parent}&=0/1.
 \end{aligned}
 \label{eq:v9-config-space-final-ledgers}
\end{equation}

\begin{theorem}[Общее finite-geometry configuration space]
Три endpoint-фазы допускают faithful orthogonal configuration carrier
$H_\Sigma$ размерности $68$ с phase algebra $\mathbb C^3$, connected
path-graph и self-adjoint block-Dirac ранга $46$. Архитектура не выбирает
финальную фазу. Канонические inclusion-edges сохраняют 63-мерный старый
подноситель и не достигают ни одной из трёх новых физических линий; кроме
того, charged increment требует family-vector edge.
\end{theorem}

\begin{nogocriterion}[Связный граф фаз не создаёт новые endpoints]
Нельзя считать ненулевые off-diagonal blocks между fibres оператором
рождения, если они являются лишь изометрическими продолжениями старых
состояний. Необходим ненулевой matrix element из исходной фазы в defect
subspace и его typed family/gauge/Real completion.
\end{nogocriterion}

\begin{protocol}\begin{enumerate}
\item число finite-geometry phases: \textbf{$3$};
\item phase algebra: \textbf{$\mathbb C^3$};
\item common configuration carrier dimension: \textbf{$68$};
\item inclusion defects: \textbf{$2$ и $1$};
\item path graph connected: \textbf{да, $\rank L_F=2$};
\item block-Dirac rank/nullity: \textbf{$46/22$};
\item target vertex выбран: \textbf{нет};
\item unreachable configuration directions: \textbf{$5$};
\item созданные physical endpoint-lines: \textbf{$0/3$};
\item configuration architecture: \textbf{$9/9$};
\item следующий гейт: \textbf{архитектура typed creation-оператора}.
\end{enumerate}\end{protocol}

Машинный сертификат сохранён в
\path{s2t_v9_endpoint_finite_geometry_configuration_space_admission_gate_results.json}.

\begin{thebibliography}{9}
\bibitem{v9-config-space-krajewski}
T.~Krajewski, \emph{Classification of Finite Spectral Triples},
arXiv:hep-th/9701081.
\bibitem{v9-config-space-quiver-action}
C.~I.~Perez-Sanchez, \emph{The Spectral Action on quivers},
arXiv:2401.03705.
\end{thebibliography}